% Generated draft: requires mathematical and source review.
\documentclass[11pt]{article}
\usepackage[margin=25mm]{geometry}
\usepackage{fontspec}
\setmainfont{Latin Modern Roman}
\usepackage{amsmath,amssymb,mathtools,unicode-math}
\setmathfont{Latin Modern Math}
\usepackage{enumitem,needspace,xurl,hyperref}
\setlength{\parindent}{0pt}
\setlength{\parskip}{5pt}
\setlength{\emergencystretch}{3em}
\newcommand{\N}{\mathbb{N}}
\newcommand{\Z}{\mathbb{Z}}
\newcommand{\Q}{\mathbb{Q}}
\newcommand{\R}{\mathbb{R}}
\newcommand{\C}{\mathbb{C}}
\newcommand{\F}{\mathbb{F}}
\newcommand{\A}{\mathbb{A}}
\newcommand{\PP}{\mathbb{P}}
\newcommand{\E}{\mathbb{E}}
\newcommand{\eps}{\varepsilon}
\newcommand{\dd}{\,\mathrm{d}}
\begin{document}
\section*{UnsolvedMath ID: 30005232}
\textbf{Source problem code:} catalogue entry
\subsection{Definitions and mathematical statement}
Let \(N\) denote the particle number and \(\ell>0\) the confining box length. The source considers the thermodynamic scaling
\[
N\longrightarrow\infty,\qquad \ell=\nu N,\qquad \nu>0,
\]
so that the macroscopic particle density is \(\nu^{-1}\).

The supplied excerpt describes conserved quantities using a matrix \(L\) and the formal expression
\[
\sum_{n=0}^{\infty}\mu_n\operatorname{tr}(L^n)
=\operatorname{tr}V(L).
\]
Here \(V\) is the potential entering the generalized Gibbs ensemble (GGE), and \(\mu_n\) are its formal coefficients. The excerpt does not define \(L\), its spectral domain, or the admissible coefficients, and gives no convergence or functional-calculus hypotheses for this expression. The term indexed by \(n=0\) is retained: although the constant function gives a trivial conservation law in the Hamiltonian description, the source emphasizes its relevance to particle density in hydrodynamics.

A separate box potential \(V_{\mathrm{box},\ell}\) is introduced to keep particle positions essentially within a box of length \(\ell\). The source notes that this auxiliary confinement generally breaks integrability and is used as a mathematical device for constructing the infinite-volume GGE. It does not specify the box potential or its boundary conditions.

Write \(Z_{\mathrm{ca},N}(\ell,V)\) for the finite-volume partition function named in the source. The displayed thermodynamic quantity is the generalized free energy per unit length,
\[
F_{\mathrm{ca}}(\nu,V)
=\lim_{N\to\infty}
\left[-\frac{1}{\nu N}\log Z_{\mathrm{ca},N}(\nu N,V)\right].
\]
This formula defines \(F_{\mathrm{ca}}\) only when the partition functions and the indicated limit are well defined; their existence is not asserted here.

\textbf{Source question, with unresolved qualifications.} Characterize precisely the confining potentials \(V\) for which the intended Calogero-type GGE has the desired infinite-volume behavior under this scaling. The excerpt identifies construction of an infinite-volume GGE, spacetime stationarity, and the displayed free energy as relevant objectives. It does not specify which of these properties constitutes the characterization, whether the free energy must be finite, or whether uniqueness or independence of the auxiliary confinement is required. The underlying Hamiltonian, matrix \(L\), phase space, reference measure, finite-volume Gibbs weight, dynamics, and mode of convergence are also absent. Consequently, the supplied record supports this qualified formulation but does not determine a fully specified necessary-and-sufficient mathematical statement.
\subsection{Short English statement}
The problem asks which confining potentials permit the intended Calogero-type generalized Gibbs ensemble to have a suitable infinite-volume limit at fixed particle density. The source also discusses stationarity and free energy per unit length. Its missing model definitions and unspecified limiting properties prevent a complete precise characterization from being stated.
\subsection{Sources}
\begin{itemize}
\item Local UnsolvedMath registry: supplied complete record with permanent ID \(30005232\), title \emph{Confining Potentials for Generalized Gibbs Ensembles}, exact source problem code \texttt{catalogue entry}. The record gives proposed year 2022 and names no proposer.
\item Supplied primary attribution: \emph{Large Scale Stochastic Dynamics} (2023), Oberwolfach Reports, pp.\ 2399--2466, DOI \texttt{10.4171/owr/2022/41}. \url{https://doi.org/10.4171/owr/2022/41}. This attribution and the extracted passage are supplied evidence; the report has not been independently consulted here.
\item Supplied literature source: \emph{Revised Enskog equation for hard rods}. \url{https://doi.org/10.1088/1742-5468/ad3ccd}.
\item Supplied literature source: \emph{Oberwolfach Report 39/2022}. \url{https://ems.press/content/serial-article-files/46978}.
\item Literature assessment dated exactly \texttt{2026-08-22}: the supplied curation labels the problem open and reports that hydrodynamic treatments describe stationary states and special integrable regimes, while no precise general class of confining potentials for the Calogero-type GGE was found. This is supplied dated curation derived from the OWR material, not an independently verified literature assessment. The record requests independent verification of time-sensitive status and cited claims.
\end{itemize}
\end{document}
